% GENERATED FILE. Edit canonical lessons, then run npm run generate:books.
% canonical-source-hash: fnv1a64:6378e5e3a79d2c6c
% canonical-lessons: ZH-C333-chunwan, ZH-C333-chuxi, ZH-C333-nianyefan, ZH-C333-shousui, ZH-C333-bainian

\chapter{New Year's Eve}
\label{ch:zh-new-year-s-eve}
\hlchaptermodality{333}

\begin{chapteropening}
\textbf{By the end of this chapter:} \emph{I can talk about New Year's Eve, the gala on TV, the reunion dinner, staying up and New Year visits.}

Complete the last lesson of chapter 333: I can talk about New Year's Eve, the gala on TV, the reunion dinner, staying up and New Year visits.
\end{chapteropening}

\section[\texorpdfstring{\zh{春晚} (Chūnwǎn)}{春晚 (Chūnwǎn)}]{\zh{春晚} (Chūnwǎn) — the Spring Festival Gala}

\label{lesson:ZH-C333-chunwan}



\begin{quote}
Before the new one: say the Mandarin for how absurd, then the Mandarin for the human world below.
\end{quote}

\subsection*{You'll want to know: \zh{春晚}}
\textbf{\zh{春晚}} — \emph{Chūnwǎn} — \textquotedblleft{}the Spring Festival Gala\textquotedblright{}.

\begin{quote}
\textbf{\zh{每年春节我们都在家看春晚}} \emph{měinián Chūnjié wǒmen dōu zài jiā kàn Chūnwǎn} — Every lunar new year we have the gala on at home.
\end{quote}

\begin{cousinweb}
\textbf{\zh{春晚}} (\emph{Chūnwǎn}) is built from \zh{春} 'spring' and \zh{晚} 'evening', a short form of the longer name of the Spring Festival evening show. Both are inherited from Middle Chinese, \emph{tsyhwin} and \emph{mjonX} (Guangyun, in Baxter's transcription), which calls the first the first of the four seasons and glosses the second 'dusk'.
\end{cousinweb}

\begin{grammarlens}[title={What you've built}]
One more word for telling what happened.
\end{grammarlens}

\subsection*{Your turn}
\begin{itemize}
\raggedright
  \item \emph{Say it:} \emph{Chūnwǎn}
  \item \emph{Say it:} \emph{Chūnwǎn}, once more
  \item \emph{Say it:} \emph{měinián Chūnjié wǒmen dōu zài jiā kàn Chūnwǎn} — Every lunar new year we have the gala on at home.
  \item \emph{Recall:} say the Mandarin for how absurd, then the Mandarin for the human world below, then say \emph{Chūnwǎn} again
\end{itemize}

\begin{tcolorbox}[breakable,colback=teal!4,colframe=teal!35!black,title={Before you move on}]
What does \zh{春晚} mean? (\textquotedblleft{}The Spring Festival Gala\textquotedblright{}.) Say it once more.
\end{tcolorbox}
\section[Chúxī]{Chúxī — lunar New Year's Eve}

\label{lesson:ZH-C333-chuxi}



\begin{quote}
Before the new one: say the Mandarin for the human world below, then the Mandarin for the Spring Festival Gala on TV.
\end{quote}

\subsection*{You'll want to know: Chúxī}
\emph{Chúxī} — \textquotedblleft{}lunar New Year's Eve\textquotedblright{}.

\begin{quote}
\emph{Chúxī de shíhou wǒmen dōu huí jiā} — On the eve of the lunar new year we all go home.
\end{quote}

Its characters come later; for now it is a word for the ear.

\begin{cousinweb}
\emph{Chúxī} is built from \emph{chú} 'to clear away' and \emph{xī} 'dusk', both inherited from Middle Chinese, \emph{drjo} and \emph{zjek} (Guangyun, in Baxter's transcription), which glosses the first 'steps' and 'to remove', and the second 'dusk'.
\end{cousinweb}

\begin{grammarlens}[title={What you've built}]
One more word for telling what happened.
\end{grammarlens}

\subsection*{Your turn}
\begin{itemize}
\raggedright
  \item \emph{Say it:} \emph{Chúxī}
  \item \emph{Say it:} \emph{Chúxī}, once more
  \item \emph{Say it:} \emph{Chúxī de shíhou wǒmen dōu huí jiā} — On the eve of the lunar new year we all go home.
  \item \emph{Recall:} say the Mandarin for the human world below, then the Mandarin for the Spring Festival Gala on TV, then say \emph{Chúxī} again
\end{itemize}

\begin{tcolorbox}[breakable,colback=teal!4,colframe=teal!35!black,title={Before you move on}]
What does Chúxī mean? (\textquotedblleft{}Lunar New Year's Eve\textquotedblright{}.) Say it once more.
\end{tcolorbox}
\section[niányèfàn]{niányèfàn — New Year's Eve dinner}

\label{lesson:ZH-C333-nianyefan}



\begin{quote}
Before the new one: say the Mandarin for the Spring Festival Gala on TV, then the Mandarin for lunar New Year's Eve.
\end{quote}

\subsection*{You'll want to know: niányèfàn}
\emph{niányèfàn} — \textquotedblleft{}New Year's Eve dinner\textquotedblright{}.

\begin{quote}
\emph{Chúxī de wǎnshang wǒmen zài nǎinai jiā chī le niányèfàn} — On the last night of the year we had the reunion dinner at my grandma's.
\end{quote}

Its characters come later; for now it is a word for the ear.

\begin{cousinweb}
\emph{niányèfàn} is built from \zh{年} 'a year', \emph{yè} 'night' and \zh{饭} 'a meal', all three inherited from Middle Chinese, \emph{nen}, \emph{yaeH} and \emph{bjonH} (Guangyun, in Baxter's transcription), which glosses the second 'evening' and says of the third that grain was first cooked into it.
\end{cousinweb}

\begin{grammarlens}[title={What you've built}]
One more word for telling what happened.
\end{grammarlens}

\subsection*{Your turn}
\begin{itemize}
\raggedright
  \item \emph{Say it:} \emph{niányèfàn}
  \item \emph{Say it:} \emph{niányèfàn}, once more
  \item \emph{Say it:} \emph{Chúxī de wǎnshang wǒmen zài nǎinai jiā chī le niányèfàn} — On the last night of the year we had the reunion dinner at my grandma's.
  \item \emph{Recall:} say the Mandarin for the Spring Festival Gala on TV, then the Mandarin for lunar New Year's Eve, then say \emph{niányèfàn} again
\end{itemize}

\begin{tcolorbox}[breakable,colback=teal!4,colframe=teal!35!black,title={Before you move on}]
What does niányèfàn mean? (\textquotedblleft{}New Year's Eve dinner\textquotedblright{}.) Say it once more.
\end{tcolorbox}
\section[shǒusuì]{shǒusuì — to keep the New Year vigil}

\label{lesson:ZH-C333-shousui}



\begin{quote}
Before the new one: say the Mandarin for lunar New Year's Eve, then the Mandarin for the New Year's Eve dinner.
\end{quote}

\subsection*{You'll want to know: shǒusuì}
\emph{shǒusuì} — \textquotedblleft{}to keep the New Year vigil\textquotedblright{}.

\begin{quote}
\emph{chī le niányèfàn yǐhòu wǒmen yìqǐ shǒusuì} — After the reunion dinner we stayed up together.
\end{quote}

Its characters come later; for now it is a word for the ear.

\begin{cousinweb}
\emph{shǒusuì} is built from \emph{shǒu} 'to guard' and \zh{岁} 'a year', both inherited from Middle Chinese, \emph{syuwX} and \emph{sjwejH} (Guangyun, in Baxter's transcription), which glosses the first 'to keep watch over'.
\end{cousinweb}

\begin{grammarlens}[title={What you've built}]
One more word for telling what happened.
\end{grammarlens}

\subsection*{Your turn}
\begin{itemize}
\raggedright
  \item \emph{Say it:} \emph{shǒusuì}
  \item \emph{Say it:} \emph{shǒusuì}, once more
  \item \emph{Say it:} \emph{chī le niányèfàn yǐhòu wǒmen yìqǐ shǒusuì} — After the reunion dinner we stayed up together.
  \item \emph{Recall:} say the Mandarin for lunar New Year's Eve, then the Mandarin for the New Year's Eve dinner, then say \emph{shǒusuì} again
\end{itemize}

\begin{tcolorbox}[breakable,colback=teal!4,colframe=teal!35!black,title={Before you move on}]
What does shǒusuì mean? (\textquotedblleft{}To keep the New Year vigil\textquotedblright{}.) Say it once more.
\end{tcolorbox}
\section[bàinián]{bàinián — to pay a New Year call}

\label{lesson:ZH-C333-bainian}



\begin{quote}
Before the new one: say the Mandarin for the New Year's Eve dinner, then the Mandarin for to keep the New Year vigil.
\end{quote}

\subsection*{You'll want to know: bàinián}
\emph{bàinián} — \textquotedblleft{}to pay a New Year call\textquotedblright{}.

\begin{quote}
\emph{dì'èr tiān wǒmen qù yéye jiā bàinián} — The next day we went to my grandpa's to wish him a happy new year.
\end{quote}

Its characters come later; for now it is a word for the ear.

\begin{cousinweb}
\emph{bàinián} is built from \emph{bài} 'to bow' and \zh{年} 'a year', both inherited from Middle Chinese, \emph{peajH} and \emph{nen} (Guangyun, in Baxter's transcription), which lists the first among the forms of a bow.
\end{cousinweb}

\begin{grammarlens}[title={What you've built}]
One more word for telling what happened.
\end{grammarlens}

\subsection*{Your turn}
\begin{itemize}
\raggedright
  \item \emph{Say it:} \emph{bàinián}
  \item \emph{Say it:} \emph{bàinián}, once more
  \item \emph{Say it:} \emph{dì'èr tiān wǒmen qù yéye jiā bàinián} — The next day we went to my grandpa's to wish him a happy new year.
  \item \emph{Recall:} say the Mandarin for the New Year's Eve dinner, then the Mandarin for to keep the New Year vigil, then say \emph{bàinián} again
\end{itemize}

\begin{tcolorbox}[breakable,colback=teal!4,colframe=teal!35!black,title={Before you move on}]
What does bàinián mean? (\textquotedblleft{}To pay a New Year call\textquotedblright{}.) Say it once more.
\end{tcolorbox}
